% Chương 2 — Cơ sở lý thuyết. Ưu tiên tài liệu trong 5 năm gần đây; không dùng MDPI, Frontiers, Hindawi.
\chapter{Cơ sở lý thuyết}\label{ch:theory}

Chương này đặt nền cho các lựa chọn phương pháp ở Chương~\ref{ch:method}. Mục~\ref{sec:concepts} định nghĩa các khái niệm dùng xuyên suốt khóa luận, gồm sự hài lòng và ``bất mãn ẩn'', thời điểm dự đoán và nhãn yếu, hai tính chất cần có của một mô hình trên dữ liệu mất cân bằng, và ngân sách can thiệp. Mục~\ref{sec:related} điểm lại năm nhóm nghiên cứu liên quan và rút ra khoảng trống mà đề tài nhắm tới. Mục~\ref{sec:theory-models} trình bày cơ sở lý thuyết của từng bước trong pipeline; mỗi phần gắn với một quyết định cụ thể ở Chương~\ref{ch:method}.

% =====================================================================
\section{Tổng quan các khái niệm liên quan}\label{sec:concepts}

\subsection{Sự hài lòng, đánh giá trực tuyến và bất mãn ẩn}\label{sec:con-satisfaction}

Theo mô hình kỳ vọng--không xác nhận (expectation--disconfirmation) của \textcite{oliver1980}, khách hình thành sự hài lòng bằng cách so sánh kết quả nhận được với kỳ vọng có từ trước khi mua. Vượt kỳ vọng thì hài lòng; hụt kỳ vọng thì bất mãn. Trong thương mại điện tử, kỳ vọng ấy có nhiều chiều, từ nhận đủ hàng, đúng sản phẩm, đúng mô tả, hàng còn nguyên vẹn cho đến nhận đúng ngày. Ngày giao dự kiến chỉ đo chiều cuối cùng. Vì vậy một đơn đến đúng hạn vẫn có thể hụt kỳ vọng ở mọi chiều còn lại. \textcite{ravula2023} còn cho thấy khách thưởng rất ít cho giao sớm nhưng phạt nặng giao trễ; nói cách khác, giao đúng hạn chủ yếu giúp tránh bị phạt chứ không bảo đảm khách hài lòng.

Sự hài lòng không quan sát được trực tiếp. Dữ liệu chỉ có điểm đánh giá từ 1 đến 5 sao mà khách tự nguyện để lại sau khi mua, và điểm này là một đại lượng đại diện với hai giới hạn. Nó chỉ tồn tại ở những khách chọn lên tiếng~\parencite{li2008,schoenmueller2020}, và nó nén một trải nghiệm nhiều chiều vào một con số duy nhất. Theo quy ước của \textcite{orman2026}, đề tài xem đánh giá 1--2 sao là không hài lòng và 3--5 sao là hài lòng. Cách chia này giữ cho lớp dương chỉ chứa phản hồi tiêu cực rõ ràng, đồng thời cho phép đối chiếu trực tiếp với nghiên cứu trước trên cùng bộ dữ liệu.

Trên nền đó, một đơn được xem là \emph{bất mãn ẩn} khi thỏa đồng thời ba điều kiện, đó là được giao thành công không muộn hơn ngày giao dự kiến (i), có đánh giá được tạo sau thời điểm giao (ii), và đánh giá đó là 1 hoặc 2 sao (iii). Gọi là ``ẩn'' vì các chỉ số vận hành quen thuộc như tỉ lệ giao trễ không ghi nhận điều gì bất thường ở đơn cho đến khi đánh giá xuất hiện, và lúc đó đã quá muộn để can thiệp. Điều kiện (ii) loại các đánh giá viết trước khi khách nhận hàng, bởi Chương~\ref{ch:results} cho thấy những đánh giá này chủ yếu phản ánh việc chưa nhận được hàng chứ không phải trải nghiệm với sản phẩm (mục~\ref{subsec:eda-lateness}).

\subsection{Thời điểm dự đoán, đơn vị phân tích và nhãn yếu}\label{sec:con-units}

Một bài toán dự báo được xác định trước hết bởi thời điểm dự đoán. Trong đề tài, thời điểm dự đoán $t_o$ của đơn $o$ là lúc đơn được giao cho khách. Đây là mốc sớm nhất mà toàn bộ thông tin giao hàng đã biết, và cũng là mốc muộn nhất mà bộ phận chăm sóc khách hàng còn có thể liên hệ trước khi khách viết đánh giá. Một đặc trưng chỉ hợp lệ khi tính được hoàn toàn từ dữ liệu có trước $t_o$; dùng thông tin không thể biết tại $t_o$ là \emph{rò rỉ dữ liệu}, được hình thức hóa ở mục~\ref{sec:th-pit}.

Đề tài làm việc với hai đơn vị phân tích. \emph{Item}, tức một sản phẩm trong đơn, là đơn vị học, vì phần lớn thông tin dự báo như danh mục, người bán, kích thước và phí vận chuyển gắn với sản phẩm. \emph{Đơn hàng} là đơn vị của nhãn, của phép chia dữ liệu và của vận hành, vì Olist chỉ ghi đánh giá theo đơn và nhân viên liên hệ khách theo đơn. Khi mỗi item thừa hưởng nhãn của đơn chứa nó, ta có một \emph{nhãn yếu}, tức nhãn đúng ở mức đơn nhưng có thể sai ở mức item, chẳng hạn khi chỉ một trong hai sản phẩm của đơn thực sự có vấn đề. Điểm của một item vì vậy là tín hiệu để ưu tiên đơn, không phải bằng chứng cho thấy item đó gây ra bất mãn. Việc nhãn nằm ở mức đơn trong khi dự đoán nằm ở mức item đưa bài toán vào khuôn khổ học đa thể hiện, được bàn ở mục~\ref{sec:th-aggregation}.

\subsection{Mất cân bằng lớp: khả năng phân biệt và độ tin cậy của xác suất}\label{sec:con-imbalance}

Trong quần thể nghiên cứu, khoảng 9\% số đơn nhận đánh giá thấp. Với tỉ lệ như vậy, cần tách hai tính chất của một mô hình xác suất. \emph{Khả năng phân biệt} (discrimination) là mức độ mô hình xếp đơn dương lên trên đơn âm; đề tài đo nó bằng PR-AUC. \emph{Độ hiệu chỉnh} (calibration) là mức độ xác suất dự đoán khớp với tần suất thực tế. Chẳng hạn, trong nhóm đơn được dự đoán 20\%, khoảng 20\% phải thực sự bị đánh giá thấp. Đề tài đo tính chất này bằng Brier score. Hai tính chất này không nhất thiết đi cùng nhau. Một mô hình có thể xếp hạng tốt mà xác suất vẫn lệch, và các kỹ thuật xử lý mất cân bằng như lấy mẫu lại hay trọng số lớp thường không cải thiện khả năng phân biệt nhưng làm xác suất bị đẩy lên một cách có hệ thống~\parencite{vandengoorbergh2022}. Đề tài cần cả hai, vì thứ tự xếp hạng quyết định đơn nào vào danh sách ưu tiên, còn xác suất giúp quản lý ước lượng khối lượng công việc. Cơ sở lý thuyết của hai loại chỉ số này nằm ở mục~\ref{sec:th-calibration} và~\ref{sec:th-metrics}.

\subsection{Ngân sách can thiệp: dự báo rủi ro và hiệu quả can thiệp}\label{sec:con-budget}

Bộ phận chăm sóc khách hàng chỉ liên hệ được một tỉ lệ $b$ số đơn mỗi kỳ; đề tài gọi tỉ lệ này là \emph{ngân sách} hay \emph{năng lực xử lý}. Câu hỏi vận hành vì vậy không phải ``đơn này có bị đánh giá thấp không'' mà là ``trong giới hạn $b$, nên chọn những đơn nào'', và thước đo tự nhiên là tỉ lệ đơn bị đánh giá thấp nằm trong nhóm được chọn, gọi là Recall@$b$ (mục~\ref{sec:th-metrics}). Cần phân biệt hai mục tiêu chọn đơn. Mục tiêu thứ nhất là \emph{tìm} các đơn có rủi ro cao; đây là bài toán dự báo. Mục tiêu thứ hai là \emph{cứu} các đơn mà việc liên hệ sẽ thay đổi được kết quả; đây là bài toán ước lượng hiệu quả can thiệp (uplift), và nó cần dữ liệu có nhóm được can thiệp và nhóm đối chứng~\parencite{ascarza2018}. Olist không ghi nhận can thiệp nào, nên đề tài chỉ giải quyết mục tiêu thứ nhất; ranh giới này được hình thức hóa ở cuối mục~\ref{sec:th-aggregation}.

% =====================================================================
\section{Các nghiên cứu liên quan}\label{sec:related}

Mục này điểm lại năm nhóm nghiên cứu mà đề tài kế thừa hoặc phản biện, gồm các mô hình dự báo sự hài lòng trên dữ liệu thương mại điện tử, bằng chứng về tác động của giao hàng đến điểm đánh giá, cơ chế tự chọn khi khách viết đánh giá, các nghiên cứu về nhắm mục tiêu can thiệp dưới ngân sách, và các nghiên cứu phương pháp về mất cân bằng lớp và rò rỉ dữ liệu. Bảng~\ref{tab:related} tóm tắt các công trình chính, và mục~\ref{sec:rel-gaps} nêu khoảng trống mà đề tài nhắm tới.

\subsection{Dự báo sự hài lòng và điểm đánh giá trong thương mại điện tử}\label{sec:rel-satisfaction}

Bộ dữ liệu Olist~\parencite{olist2018} là nền tảng thực nghiệm phổ biến cho bài toán phân loại sự hài lòng, vì nó nối được đánh giá của khách với thông tin đơn, sản phẩm, người bán và vận chuyển. Công trình gần nhất với đề tài là của \textcite{orman2026}, dùng cùng quy ước nhãn (1--2 sao là không hài lòng, 3--5 sao là hài lòng) trên 108\,449 quan sát. Tác giả loại các biến văn bản đánh giá để giảm rò rỉ, so sánh tám thuật toán học máy truyền thống (trong đó có Random Forest, CatBoost, XGBoost, LightGBM) với năm kiến trúc học sâu, chia dữ liệu ngẫu nhiên có phân tầng theo tỉ lệ 80:20 và chỉ áp dụng SMOTE trên tập huấn luyện. Mô hình GRU-LSTM tốt nhất đạt accuracy 86,12\%, nhưng với lớp không hài lòng chỉ đạt recall 14,35\% và F1 24,14\%; balanced accuracy là 56,76\%. Chính tác giả kết luận rằng hiệu năng tổng thể cao có thể che giấu khả năng nhận diện lớp thiểu số rất yếu.

Hai điểm thiết kế của nghiên cứu này khác với bối cảnh vận hành mà đề tài hướng tới. Thứ nhất, phép chia ngẫu nhiên cho phép mô hình học từ các đơn xảy ra sau đơn đang được đánh giá, điều không thể có khi triển khai. Thứ hai, tập đặc trưng có thuộc tính hoàn tất đơn và trạng thái đơn, và mẫu bao gồm cả các đơn giao trễ, vốn là nhóm mà riêng tín hiệu trễ hạn đã đủ để cảnh báo. Cả hai lựa chọn đều có xu hướng làm hiệu năng đo được cao hơn mức mà một hệ thống triển khai thực tế có thể đạt.

\textcite{zaghloul2024} so sánh các mô hình truyền thống như Random Forest và SVM với các mạng nơ-ron trên hơn 100 nghìn đơn hàng trực tuyến và báo cáo Random Forest đạt accuracy 92\%, cao hơn các mô hình học sâu; thời gian giao và độ chính xác của đơn nằm trong nhóm yếu tố quan trọng nhất. Kết quả này cùng chiều với bằng chứng tổng quát hơn về dữ liệu bảng. \textcite{shwartzziv2022} cho thấy XGBoost vượt các mô hình học sâu chuyên cho dữ liệu bảng trên nhiều bộ dữ liệu, kể cả các bộ do chính tác giả của các mô hình học sâu đó chọn, và cần ít công tinh chỉnh hơn. Tuy nhiên, accuracy không cho biết mô hình bắt được bao nhiêu đơn bị đánh giá thấp. Khi lớp đó chỉ chiếm khoảng một phần mười, một mô hình gán mọi đơn là hài lòng cũng đạt accuracy xấp xỉ 90\%. Các kết quả trên vì thế ủng hộ việc chọn mô hình cây làm điểm khởi đầu, nhưng chưa trả lời câu hỏi vận hành của đề tài.

\subsection{Hiệu quả giao hàng và đánh giá trực tuyến}\label{sec:rel-delivery}

Ba nghiên cứu thực nghiệm giải thích vì sao giao hàng chi phối điểm đánh giá và vì sao ``đúng hạn'' vẫn chưa đủ. \textcite{ravula2023} ước lượng mô hình logit thứ bậc trong khung Bayes và thấy khách chấm điểm thấp hơn nhiều cho đơn giao trễ so với đơn đúng hạn, trong khi đơn giao sớm chỉ được chấm cao hơn một chút. \textcite{deshpande2023} dùng hơn 15 triệu đơn trên Tmall và mạng logistics Cainiao, kết hợp mô hình hồi quy thứ bậc kiểu Heckman để giải thích đồng thời điểm đánh giá logistics và khả năng khách có để lại đánh giá, rồi nối điểm logistics với xác suất mua và doanh số của người bán. Nghiên cứu này cho thấy chất lượng giao hàng tác động đến doanh thu thông qua đánh giá, nên một đánh giá thấp có hệ quả kinh tế vượt ra ngoài một đơn đơn lẻ.

\textcite{brylla2022} chỉ ra một ma sát mà chỉ số đúng hạn không ghi nhận. Qua ba thí nghiệm và một nghiên cứu thực địa, các tác giả thấy khi đơn bị tách thành nhiều lần giao, giao nhanh hơn không cải thiện mức hoàn tất đơn, sự hài lòng hay truyền miệng; tác động tiêu cực đi qua chi phí phiền toái và cảm nhận về tác hại môi trường, và dữ liệu bán hàng thực tế cho thấy giao tách làm giảm số lần mua lại. Đặt cạnh nhau, ba kết quả này gợi ý rằng trong nhóm đơn đúng hạn vẫn còn những nguồn bất mãn gắn với thời gian chờ tuyệt đối, cấu trúc đơn và chất lượng vận hành của người bán. Olist không ghi thời điểm giao riêng cho từng item, nên việc tách kiện chỉ quan sát được gián tiếp qua số người bán và số sản phẩm trong đơn (Bảng~\ref{tab:features}).

\subsection{Thiên lệch tự chọn trong đánh giá trực tuyến}\label{sec:rel-selection}

Đánh giá trực tuyến không phải mẫu ngẫu nhiên của trải nghiệm khách hàng. \textcite{li2008} mô hình hóa trường hợp người mua sớm có sở thích khác người mua muộn, khiến chuỗi điểm đánh giá của một sản phẩm thay đổi có hệ thống theo thời gian; kiểm định trên đánh giá sách của Amazon, các tác giả cho thấy thiên lệch tự chọn nếu không được hiệu chỉnh sẽ làm giảm thặng dư người tiêu dùng. \textcite{schoenmueller2020} phân tích hơn 280 triệu đánh giá trên 25 nền tảng và xác nhận hiện tượng ``tự chọn theo cực'', trong đó người có trải nghiệm rất tốt hoặc rất tệ có xu hướng viết đánh giá nhiều hơn, làm phân phối điểm lệch về hai đầu và giảm lượng thông tin mà đánh giá mang lại. Mô hình Heckman của \textcite{deshpande2023} là một ví dụ xử lý trực tiếp quá trình chọn này bằng cách mô hình hóa riêng xác suất khách có đánh giá.

Hệ quả cho đề tài là nhãn dương chỉ có nghĩa ``đơn nhận đánh giá 1--2 sao trong số các đơn có đánh giá''. Khách không viết đánh giá chưa chắc đã hài lòng, và các yếu tố dự báo việc khách lên tiếng có thể khác các yếu tố dự báo việc khách chấm điểm thấp. Đề tài không mô hình hóa riêng quá trình có đánh giá; giới hạn này được ghi nhận khi phát biểu bài toán ở mục~\ref{sec:pipeline}.

\subsection{Dự báo rủi ro và ưu tiên can thiệp dưới ngân sách}\label{sec:rel-targeting}

Nghiên cứu về giữ chân khách hàng cảnh báo rằng xếp hạng theo rủi ro và xếp hạng theo hiệu quả can thiệp là hai việc khác nhau. Kết hợp hai thí nghiệm thực địa với học máy, \textcite{ascarza2018} cho thấy khách có rủi ro rời bỏ cao nhất không nhất thiết là đối tượng nên nhắm tới, và nhắm theo mức nhạy cảm với can thiệp hiệu quả hơn đáng kể so với nhắm theo rủi ro. \textcite{devriendt2021} đi tiếp theo hướng này khi đề xuất thước đo lợi nhuận tối đa cho mô hình uplift và, trong một nghiên cứu tình huống ở ngành tài chính, so sánh trực tiếp mô hình uplift với mô hình dự báo rời bỏ. Cả hai đều cần dữ liệu có nhóm được can thiệp và nhóm đối chứng, thứ mà Olist không có.

Khi chưa có dữ liệu can thiệp, câu hỏi khả thi là chọn đơn nào trong giới hạn năng lực xử lý. \textcite{gomezvargas2025} phát biểu việc chọn khách cho chiến dịch giữ chân như một bài toán tối thiểu hóa hối tiếc theo khung predict-and-optimize, dùng giá trị vòng đời của từng khách thay vì giá trị trung bình, và đạt lợi nhuận trung bình tốt nhất trên 12 bộ dữ liệu rời bỏ. Một kết quả thận trọng hơn đến từ \textcite{li2026valuesorting}, một bản tiền ấn phẩm so sánh ưu tiên lô hàng bằng học máy với việc chỉ sắp xếp theo giá trị trên ba bộ dữ liệu, trong đó có Olist. Ở ngân sách xem xét 10\%, mô hình kém luật sắp theo giá trị 5,5 điểm phần trăm trên SCMS và 4,9 điểm trên Olist, nhưng hơn 10,1 điểm trên DataCo. Tác giả gắn chênh lệch này với việc mức độ trễ có học được từ dữ liệu hay không và với độ lệch hiệu chỉnh, rồi khuyến nghị giữ luật sắp theo giá trị làm mốc so sánh thường trực và chỉ triển khai học máy sau khi đã qua đánh giá rolling-origin. Bài toán của đề tài khác ở chỗ nhãn là đánh giá thấp chứ không phải mức độ trễ, nhưng bài học về mốc so sánh đơn giản vẫn áp dụng được.

\subsection{Mất cân bằng lớp, rò rỉ dữ liệu và độ tin cậy của kết quả}\label{sec:rel-imbalance}

Hai nghiên cứu phương pháp định hướng cách xử lý mất cân bằng trong đề tài. Bằng mô phỏng Monte Carlo và một nghiên cứu tình huống chẩn đoán ung thư buồng trứng, \textcite{vandengoorbergh2022} cho thấy undersampling, oversampling và SMOTE đều tạo ra mô hình hồi quy logistic có xác suất lớp thiểu số bị ước lượng quá cao, trong khi diện tích dưới đường ROC không tăng; sự cân bằng giữa độ nhạy và độ đặc hiệu mà các phương pháp này mang lại cũng đạt được chỉ bằng cách dịch ngưỡng xác suất. \textcite{elor2022smote} đi đến kết luận tương tự cho các bộ phân loại mạnh. Theo nhóm tác giả, cân bằng dữ liệu có ích cho bộ phân loại yếu nhưng hầu như không cải thiện các mô hình boosting như LightGBM, XGBoost và CatBoost khi ngưỡng quyết định đã được tối ưu.

Rò rỉ dữ liệu là mối đe dọa thứ hai đối với độ tin cậy của kết quả. Khảo sát của \textcite{kapoor2023} tìm thấy rò rỉ trong 294 bài báo thuộc 17 lĩnh vực ứng dụng học máy và đề xuất phân loại tám dạng rò rỉ, từ lỗi sơ đẳng như tiền xử lý trên cả tập huấn luyện và kiểm tra đến rò rỉ theo thời gian và sự phụ thuộc giữa các tập. Khi sửa các lỗi này trong bài toán dự báo nội chiến, các mô hình học máy phức tạp không còn vượt trội đáng kể so với hồi quy logistic. Kết quả này giải thích vì sao Chương~\ref{ch:method} ưu tiên kiểm soát rò rỉ trước khi so sánh mô hình.

\begin{table}[htbp]\centering\footnotesize
\caption{Tóm tắt các nghiên cứu liên quan}\label{tab:related}
\begin{tabular}{p{2.5cm}p{3.0cm}p{3.4cm}p{5.3cm}}
\toprule
Nghiên cứu & Dữ liệu & Phương pháp & Kết quả chính và hạn chế đối với đề tài \\
\midrule
\textcite{orman2026} & Olist, 108\,449 quan sát, nhãn 1--2 sao & 8 mô hình ML, 5 mô hình học sâu; chia ngẫu nhiên 80:20; SMOTE & Accuracy 86,12\% nhưng recall lớp không hài lòng 14,35\%; không đánh giá theo thời gian, có cả đơn giao trễ \\
\textcite{zaghloul2024} & Hơn 100 nghìn đơn trực tuyến & RF, SVM so với mạng nơ-ron & RF đạt accuracy 92\%; chỉ số tổng thể không phản ánh khả năng bắt lớp thiểu số \\
\textcite{ravula2023} & Đánh giá đơn hàng trực tuyến & Logit thứ bậc Bayes & Giao trễ bị phạt nặng, giao sớm chỉ được thưởng nhẹ; mô hình giải thích, không dự báo \\
\textcite{deshpande2023} & Hơn 15 triệu đơn Tmall/Cainiao & Hồi quy thứ bậc Heckman, mô hình lựa chọn & Điểm logistics liên hệ với xác suất mua và doanh số; có xử lý việc khách có đánh giá hay không \\
\textcite{brylla2022} & 3 thí nghiệm, 1 nghiên cứu thực địa & Thí nghiệm, phân tích trung gian & Giao tách triệt tiêu lợi ích của giao nhanh, giảm mua lại; Olist không quan sát trực tiếp giao tách \\
\textcite{schoenmueller2020} & Hơn 280 triệu đánh giá, 25 nền tảng & Phân tích đa nền tảng, thí nghiệm, khảo sát & Tự chọn theo cực làm đánh giá kém đại diện; nhãn chỉ phản ánh bất mãn được quan sát \\
\textcite{ascarza2018} & 2 thí nghiệm thực địa & Học máy trên dữ liệu thí nghiệm & Khách rủi ro cao nhất không phải đối tượng can thiệp tốt nhất; cần dữ liệu can thiệp \\
\textcite{gomezvargas2025} & 12 bộ dữ liệu rời bỏ & Predict-and-optimize, giá trị vòng đời cá nhân & Lợi nhuận trung bình tốt nhất; cần giả định về chi phí và phản ứng với can thiệp \\
\textcite{li2026valuesorting} & SCMS, DataCo, Olist & So sánh ML với sắp theo giá trị ở ngân sách 10\% & ML thua trên 2/3 bộ dữ liệu, trong đó có Olist; bản tiền ấn phẩm \\
\textcite{vandengoorbergh2022} & Mô phỏng, dữ liệu ung thư buồng trứng & Logistic với 4 cách xử lý mất cân bằng & Lấy mẫu lại làm xác suất bị đẩy lên, AUC không tăng \\
\textcite{elor2022smote} & Nhiều bộ dữ liệu mất cân bằng & So sánh cân bằng dữ liệu với tối ưu ngưỡng & SMOTE không giúp các mô hình boosting mạnh \\
\textcite{kapoor2023} & 294 bài báo thuộc 17 lĩnh vực & Khảo sát, phân loại 8 dạng rò rỉ & Sửa rò rỉ xóa lợi thế của mô hình phức tạp; cần kiểm soát rò rỉ trước khi so sánh \\
\bottomrule
\end{tabular}
\end{table}

\subsection{Khoảng trống nghiên cứu}\label{sec:rel-gaps}

Tổng hợp lại, các nghiên cứu trên để lại sáu khoảng trống sau.
\begin{enumerate}
  \item \textbf{Quần thể.} Phần lớn công trình dự báo hoặc giải thích đánh giá trên toàn bộ đơn, nơi giao trễ là tín hiệu mạnh nhất. Nhóm đơn giao đúng hoặc sớm hạn mà vẫn bị chấm 1--2 sao chưa được mô hình hóa như một đối tượng riêng.
  \item \textbf{Thời điểm dự báo.} Các mô hình hài lòng hiện có dùng phép chia ngẫu nhiên và không cố định một mốc mà tại đó mọi đặc trưng đều đã biết, nên chưa cho biết hiệu năng khi triển khai.
  \item \textbf{Chỉ số đánh giá.} Các nghiên cứu này báo cáo accuracy hoặc F1 tổng thể và dùng SMOTE, trong khi bằng chứng phương pháp cho thấy lấy mẫu lại làm hỏng xác suất. Chưa có đánh giá kết hợp PR-AUC, độ hiệu chỉnh và phân tích theo ngân sách cho nhóm đơn đúng hạn.
  \item \textbf{Nhãn.} Nhãn đánh giá chịu thiên lệch tự chọn nhưng thường được dùng như sự thật nền.
  \item \textbf{Mốc so sánh.} Một mô hình có chỉ số phân biệt tốt chưa chắc ưu tiên tốt hơn luật đơn giản, như kết quả trên Olist của \textcite{li2026valuesorting} cho thấy.
  \item \textbf{Can thiệp.} Dữ liệu quan sát không có can thiệp nên không thể biết đơn nào ``cứu'' được.
\end{enumerate}

Đề tài lấp trực tiếp ba khoảng trống đầu. Quần thể được giới hạn ở đơn giao đúng hoặc sớm hạn, mọi đặc trưng được dựng tại thời điểm giao và kiểm định bằng rolling-origin theo thời gian, còn việc đánh giá dựa trên PR-AUC, Brier score và Recall@5\%, với trọng số lớp thay cho lấy mẫu lại. Khoảng trống thứ năm được xử lý một phần, vì Recall@5\% mới chỉ được so với mốc chọn ngẫu nhiên, còn phép so sánh với luật sắp theo giá trị là hướng cần bổ sung (mục~\ref{sec:eval}). Hai khoảng trống còn lại, thiên lệch tự chọn và hiệu quả nhân quả của can thiệp, nằm ngoài phạm vi khóa luận. Chúng đóng vai trò ranh giới diễn giải, nghĩa là mọi kết luận chỉ dừng ở mức dự báo và ưu tiên, không suy ra tác động của việc liên hệ khách hàng.

% =====================================================================
\section{Lý thuyết và mô hình áp dụng}\label{sec:theory-models}

Mục này trình bày nền tảng lý thuyết của các lựa chọn phương pháp ở Chương~\ref{ch:method}, theo thứ tự của pipeline, lần lượt là ràng buộc thông tin tại thời điểm dự báo, cách ước lượng hiệu năng theo thời gian, các họ mô hình cây, hiệu chỉnh xác suất, chỉ số đánh giá cho lớp thiểu số, và cách điểm item được chuyển thành danh sách ưu tiên mức đơn.

\subsection{Dự báo tại thời điểm giao hàng và rò rỉ dữ liệu}\label{sec:th-pit}

\textcite{kaufman2012} định nghĩa rò rỉ là việc đưa vào quá trình học những thông tin về biến mục tiêu mà tại thời điểm dự đoán không thể biết một cách hợp lệ, và lập luận rằng cách phòng tránh căn bản là gắn mỗi quan sát và mỗi nguồn thông tin với thời điểm nó xuất hiện. Áp dụng vào đề tài, với đơn $o$ có thời điểm dự đoán $t_o$, mỗi đặc trưng của item $i \in o$ phải là hàm của tập dữ liệu $\mathcal{D}_{<t_o}$ gồm các sự kiện đã xảy ra trước $t_o$, tức là
\begin{equation}\label{eq:th-pit}
  \mathbf{x}_i = \phi\left(\mathcal{D}_{<t_o}\right), \qquad i \in o .
\end{equation}
Đặc trưng lịch sử đòi hỏi ràng buộc chặt hơn vì nó tổng hợp nhãn của các đơn khác. Gọi $\tau_{o'}$ là thời điểm tạo đánh giá của đơn $o'$; tập đơn được phép đóng góp vào lịch sử của người bán $s$ tại thời điểm $t$ là
\begin{equation}\label{eq:th-history}
  \mathcal{H}_s(t) = \left\{ o' : \mathrm{seller}(o') = s,\; t_{o'} < t,\; \tau_{o'} < t \right\},
\end{equation}
với cận trên mở. Điều kiện $\tau_{o'} < t$ loại các đơn đã giao trước $t$ nhưng chưa có đánh giá tại $t$, vì nhãn của chúng chưa tồn tại tại thời điểm dự đoán. Khi $|\mathcal{H}_s(t)|$ nhỏ, trung bình lịch sử có phương sai lớn; đó là lý do Chương~\ref{ch:method} để trống giá trị và bật cờ hỗ trợ thấp khi số quan sát dưới ngưỡng.

Ràng buộc~\eqref{eq:th-pit} áp dụng cả cho các bước tiền xử lý. Trung vị dùng để điền giá trị thiếu hay danh sách mã one-hot cũng là thống kê của dữ liệu, nên nếu được tính trên toàn bộ mẫu, chúng mang thông tin của giai đoạn kiểm định vào mô hình. Phân loại của \textcite{kapoor2023} xếp lỗi này vào nhóm tiền xử lý trên cả tập huấn luyện và kiểm tra, tách biệt với rò rỉ theo thời gian và rò rỉ do các quan sát phụ thuộc nhau nằm ở hai phía của phép chia. Dạng phụ thuộc cuối cùng xuất hiện tự nhiên trong dữ liệu Olist, vì các item cùng đơn chia sẻ một nhãn, và được xử lý bằng cách đặt mọi item của một đơn vào cùng một phân vùng.

\subsection{Đánh giá theo thời gian và rolling-origin}\label{sec:th-temporal}

Kiểm định chéo $k$-fold ngẫu nhiên dựa trên giả định các quan sát độc lập và cùng phân phối. Dữ liệu Olist vi phạm giả định này theo hai cách. Tỉ lệ dương thay đổi giữa các giai đoạn (Bảng~\ref{tab:split}), và đặc trưng lịch sử của một đơn phụ thuộc vào các đơn trước nó. \textcite{hyndman2021} mô tả kiểm định chéo cho chuỗi thời gian dưới dạng đánh giá trên \emph{gốc dự báo cuộn} (rolling forecasting origin), trong đó mỗi tập kiểm định chỉ đi sau dữ liệu huấn luyện của nó và sai số được lấy trung bình qua các tập kiểm định. Đề tài không dự báo một chuỗi thời gian, nhưng dùng cùng nguyên tắc cho dòng sự kiện đơn hàng. Với các mốc cắt $c_1 < c_2 < \dots < c_{K+1}$, fold $k$ có
\begin{equation}\label{eq:th-rolling}
  \mathcal{T}_k = \{ o : t_o < c_k \}, \qquad
  \mathcal{V}_k = \{ o : c_k \le t_o < c_{k+1} \}, \qquad
  \bar{M} = \frac{1}{K} \sum_{k=1}^{K} M\!\left(f_{\mathcal{T}_k}; \mathcal{V}_k\right),
\end{equation}
trong đó $f_{\mathcal{T}_k}$ là mô hình học trên $\mathcal{T}_k$ và $M$ là chỉ số đánh giá. Cửa sổ huấn luyện mở rộng dần, mô phỏng việc huấn luyện lại định kỳ khi có thêm dữ liệu.

Lựa chọn này có cơ sở thực nghiệm. So sánh các phương pháp ước lượng hiệu năng trên 174 chuỗi thực và ba chuỗi mô phỏng, \textcite{cerqueira2020} thấy kiểm định chéo theo khối dùng được cho chuỗi dừng, nhưng với chuỗi không dừng, các phương pháp ngoài mẫu, đặc biệt holdout lặp lại trên nhiều giai đoạn kiểm định, cho ước lượng chính xác nhất. Thiết kế ba fold ở Hình~\ref{fig:temporal-folds} thuộc nhóm phương pháp ngoài mẫu này. Cái giá phải trả là số fold nhỏ, khiến trung bình $\bar M$ có phương sai lớn, trong khi fold đầu học trên ít dữ liệu nhất.

\subsection{Rừng ngẫu nhiên và gradient boosting}\label{sec:th-trees}

Rừng ngẫu nhiên~\parencite{breiman2001} là tập hợp $B$ cây quyết định, mỗi cây học trên một mẫu bootstrap của dữ liệu và chỉ xét một tập con ngẫu nhiên các đặc trưng tại mỗi lần tách nút. Xác suất dự đoán là trung bình của các cây
\begin{equation}\label{eq:th-rf}
  \hat{p}_{\mathrm{RF}}(\mathbf{x}) = \frac{1}{B} \sum_{b=1}^{B} T_b(\mathbf{x}),
\end{equation}
với $T_b(\mathbf{x})$ là tỉ lệ dương trong lá chứa $\mathbf{x}$ của cây $b$. \textcite{breiman2001} chứng minh sai số tổng quát hóa hội tụ khi số cây tăng, và giới hạn của nó phụ thuộc vào độ mạnh của từng cây và độ tương quan giữa các cây; việc chọn ngẫu nhiên đặc trưng nhằm giảm tương quan đó. Các ràng buộc độ sâu tối đa và số mẫu tối thiểu mỗi lá ở Chương~\ref{ch:method} kiểm soát độ phức tạp của từng cây. Trọng số lớp tính lại trên từng mẫu bootstrap làm thay đổi tỉ lệ dương hiệu dụng trong các lá, nên $\hat p_{\mathrm{RF}}$ không còn là ước lượng không chệch của xác suất thực; mục~\ref{sec:th-calibration} xử lý hệ quả này.

Gradient boosting xây mô hình cộng tính theo từng bước. Với hàm mất mát log-loss $L$ trên thang log-odds, bước $m$ thêm một cây $h_m$ khớp với gradient âm của $L$ tại mô hình hiện tại
\begin{equation}\label{eq:th-gb}
  F_m(\mathbf{x}) = F_{m-1}(\mathbf{x}) + \eta\, h_m(\mathbf{x}), \qquad
  h_m \approx \arg\min_h \sum_{j} \left( -\frac{\partial L(y_j, F)}{\partial F}\Big|_{F=F_{m-1}(\mathbf{x}_j)} - h(\mathbf{x}_j) \right)^2,
\end{equation}
với $\eta$ là tốc độ học. \textcite{prokhorenkova2018} chỉ ra rằng triển khai thông thường bị \emph{dịch chuyển dự đoán} (prediction shift), vì gradient tại mỗi bước được tính trên chính các mẫu mà mô hình đã học, nên phân phối của $F_{m-1}(\mathbf{x}_j)$ trên mẫu huấn luyện khác với trên mẫu mới; đây là một dạng rò rỉ nhãn bên trong thuật toán. Vấn đề tương tự xảy ra khi mã hóa biến phân loại bằng thống kê nhãn (target statistics) tính trên toàn bộ mẫu. CatBoost giải quyết cả hai bằng một hoán vị ngẫu nhiên $\sigma$ của dữ liệu. \emph{Ordered boosting} tính phần dư của mẫu $\sigma_k$ bằng mô hình chỉ học từ các mẫu đứng trước nó, và \emph{ordered target statistics} mã hóa giá trị phân loại của mẫu đó cũng chỉ từ các mẫu đứng trước, theo công thức
\begin{equation}\label{eq:th-ordered-ts}
  \hat{x}_{\sigma_k} =
  \frac{\sum_{j<k} \mathbf{1}\!\left[x_{\sigma_j} = x_{\sigma_k}\right] y_{\sigma_j} + a\,P}
       {\sum_{j<k} \mathbf{1}\!\left[x_{\sigma_j} = x_{\sigma_k}\right] + a},
\end{equation}
trong đó $P$ là giá trị tiên nghiệm và $a > 0$ là trọng số của nó. Công thức~\eqref{eq:th-ordered-ts} có cùng logic với~\eqref{eq:th-history} ở chỗ mỗi quan sát chỉ được nhìn thấy ``quá khứ'' của nó. Điểm khác là thứ tự trong CatBoost là hoán vị nhân tạo, còn thứ tự trong~\eqref{eq:th-history} là thời gian thực; vì vậy cơ chế của CatBoost chống rò rỉ bên trong tập huấn luyện nhưng không thay thế cho việc chia dữ liệu theo thời gian. CatBoost cũng dùng cây đối xứng (oblivious trees), trong đó mọi nút cùng độ sâu dùng chung một điều kiện tách, làm mô hình ít bị quá khớp hơn và dự đoán nhanh hơn.

\subsection{Hiệu chỉnh xác suất}\label{sec:th-calibration}

Một bộ phân loại được gọi là hiệu chỉnh tốt khi, trong nhóm quan sát nhận điểm $\hat p$, tỉ lệ dương thực tế xấp xỉ $\hat p$. \textcite{niculescumizil2005} so sánh mười thuật toán học có giám sát và thấy các phương pháp biên cực đại như boosted trees đẩy xác suất ra xa 0 và 1, tạo méo dạng chữ S, trong khi naive Bayes đẩy xác suất về hai đầu; sau khi hiệu chỉnh, boosted trees, rừng ngẫu nhiên và SVM cho xác suất tốt nhất. Hai phương pháp hiệu chỉnh được dùng trong đề tài đều học một ánh xạ từ điểm thô $s$ sang xác suất trên một tập dữ liệu không dùng để huấn luyện mô hình gốc.

\emph{Platt scaling}~\parencite{platt1999} giả định méo có dạng logistic và khớp hai tham số $A$ và $B$ bằng hợp lý cực đại cho hàm
\begin{equation}\label{eq:th-platt}
  \hat{p} = \frac{1}{1 + \exp(A s + B)} .
\end{equation}
\emph{Hồi quy isotonic}~\parencite{zadrozny2002} không giả định dạng hàm mà tìm hàm không giảm $g$ cực tiểu hóa sai số bình phương
\begin{equation}\label{eq:th-isotonic}
  \hat{g} = \arg\min_{g \text{ không giảm}} \sum_{j=1}^{n} \left( y_j - g(s_j) \right)^2 ,
\end{equation}
giải bằng thuật toán pool-adjacent-violators, cho kết quả là hàm bậc thang. Đổi lại sự linh hoạt, isotonic cần nhiều dữ liệu hiệu chỉnh hơn mới phát huy tác dụng và dễ quá khớp khi tập hiệu chỉnh nhỏ~\parencite{niculescumizil2005}. Đây là lý do Chương~\ref{ch:method} fit hai phương pháp trên một phân vùng và so sánh chúng bằng Brier score trên một phân vùng khác.

Cả hai ánh xạ đều đơn điệu, nên không đảo thứ tự xếp hạng. Platt scaling với $A<0$ là hàm tăng ngặt và giữ nguyên hoàn toàn thứ tự; isotonic là hàm bậc thang nên có thể gộp các điểm khác nhau vào cùng một giá trị, tạo ra các cặp hòa và làm PR-AUC thay đổi nhẹ. Brier score phản ánh cả hai mặt, vì nó giảm khi xác suất gần tỉ lệ thực tế hơn và cả khi mô hình tách hai lớp tốt hơn; do đó nó dùng được làm tiêu chí so sánh phương pháp hiệu chỉnh khi thứ tự xếp hạng gần như không đổi. Một mốc tham chiếu hữu ích là dự đoán hằng số bằng tỉ lệ dương $\pi$, có Brier score bằng $\pi(1-\pi)$; một mô hình có Brier cao hơn mốc này cho xác suất kém tin cậy hơn cả việc không dùng đặc trưng nào.

\subsection{Chỉ số đánh giá cho lớp thiểu số}\label{sec:th-metrics}

Gọi $\pi$ là tỉ lệ dương. Một bộ phân loại gán mọi đơn vào lớp âm có accuracy $1-\pi$, khoảng 0,9 trong quần thể của đề tài, nhưng không phát hiện được đơn nào bị đánh giá thấp. ROC-AUC không phụ thuộc vào $\pi$, nhưng chính vì vậy nó che đi chi phí của các dự đoán dương sai. Tỉ lệ dương giả được chia cho tổng số đơn âm, vốn nhiều gấp khoảng mười lần số đơn dương, nên hàng trăm dương giả chỉ làm đường ROC dịch đi rất ít. Precision thì chia cho số đơn được gắn cờ, nên phản ánh trực tiếp các dương giả đó. \textcite{saito2015} cho thấy trên dữ liệu mất cân bằng, đường ROC có thể trông tốt trong khi độ chính xác của các dự đoán dương rất thấp, và đường precision--recall phản ánh điều đó trung thực hơn. Với ngưỡng thứ $n$ theo thứ tự điểm giảm dần, precision $P_n$ và recall $R_n$ được định nghĩa thông thường; average precision (AP), ước lượng của PR-AUC dùng trong đề tài, là
\begin{equation}\label{eq:th-ap}
  \mathrm{AP} = \sum_{n} \left( R_n - R_{n-1} \right) P_n .
\end{equation}
Một mô hình xếp hạng ngẫu nhiên có precision kỳ vọng bằng $\pi$ ở mọi mức recall, nên AP của nó xấp xỉ $\pi$. Mốc này thay đổi theo fold vì $\pi$ thay đổi, và nó cũng khác nhau giữa mức item và mức đơn vì tỉ lệ dương ở hai mức khác nhau; vì vậy PR-AUC chỉ nên được đọc cùng với tỉ lệ dương tương ứng.

AP tóm tắt toàn bộ đường cong, còn quyết định vận hành chỉ quan tâm một điểm trên đường đó. Với ngân sách $b$ trên $N$ đơn, đặt $K = \lceil bN \rceil$ và $\mathrm{Top}_K$ là $K$ đơn có điểm cao nhất; khi đó
\begin{equation}\label{eq:th-recallk}
  \mathrm{Recall@}b = \frac{\sum_{o \in \mathrm{Top}_K} y_o}{\sum_{o} y_o}, \qquad
  \mathrm{Precision@}b = \frac{1}{K} \sum_{o \in \mathrm{Top}_K} y_o .
\end{equation}
Chọn ngẫu nhiên $K$ đơn cho $\mathrm{E}[\mathrm{Recall@}b] = b$, nên tỉ số $\mathrm{Recall@}b / b$ là hệ số nâng (lift) so với không dùng mô hình. Recall@$b$ chỉ phụ thuộc vào thứ tự xếp hạng, không phụ thuộc giá trị tuyệt đối của xác suất; vì vậy nó bổ sung chứ không thay thế Brier score.

\subsection{Gộp điểm item lên đơn và ưu tiên theo năng lực xử lý}\label{sec:th-aggregation}

Cấu trúc nhãn của bài toán tương ứng với học đa thể hiện (multiple-instance learning). Mỗi đơn là một ``túi'' gồm nhiều item, nhãn chỉ có ở mức túi, và giả định chuẩn của \textcite{dietterich1997} là một túi dương khi có ít nhất một thể hiện dương. Dưới giả định đó, điểm tự nhiên nhất của túi là điểm lớn nhất của các thể hiện, đúng với luật max-v1 ở~\eqref{eq:maxv1}. Giả định này chỉ gần đúng với Olist, bởi khách có thể chấm thấp vì một sản phẩm lỗi, nhưng cũng có thể vì trải nghiệm chung của cả đơn, và nhãn item kế thừa từ đơn là nhãn yếu (mục~\ref{sec:con-units}).

Luật max có một thiên lệch có thể mô tả định lượng. Nếu điểm của các item trong đơn độc lập với hàm phân phối $F$, thì
\begin{equation}\label{eq:th-maxbias}
  \Pr\left( \hat{r}_o \le t \right) = F(t)^{n_o},
\end{equation}
với $n_o$ là số item, nên phân phối của $\hat r_o$ dịch dần về phía điểm cao khi $n_o$ tăng, kể cả khi từng item không rủi ro hơn. Hai hệ quả theo sau. Đơn nhiều sản phẩm có xu hướng được xếp cao hơn; và kể cả khi xác suất item đã hiệu chỉnh tốt, $\hat r_o$ không nhất thiết là xác suất hiệu chỉnh tốt ở mức đơn. Vì vậy Brier score mức đơn~\eqref{eq:brier} cần được báo cáo riêng thay vì suy ra từ chất lượng hiệu chỉnh mức item. Mục~\ref{sec:analysis-eda} đối chiếu thiên lệch này với tỉ lệ đánh giá thấp thực tế theo số item trong đơn.

Bước cuối biến điểm thành quyết định. Bộ phận chăm sóc khách hàng chỉ xử lý được $K$ đơn mỗi kỳ, nên bài toán là chọn tập $S$ với $|S| \le K$ để cực đại hóa số đơn bất mãn được tiếp cận. Nếu $\hat r_o$ là ước lượng của $\Pr(y_o = 1)$, kỳ vọng số đơn dương trong $S$ là $\sum_{o \in S} \hat r_o$, và lời giải tối ưu là chọn $K$ đơn có điểm cao nhất. Lời giải này chỉ dùng đến thứ tự của $\hat r_o$, nên vẫn đúng khi điểm đơn chưa hiệu chỉnh tốt, miễn là nó tăng đơn điệu theo xác suất thực; đây là cơ sở của các ngưỡng top 5\% và top 20\% ở Chương~\ref{ch:method}. Luật này tối ưu số đơn bất mãn được \emph{tìm thấy}, không tối ưu số đơn được \emph{cứu}. Mục tiêu thứ hai đòi hỏi xếp hạng theo hiệu ứng can thiệp
\begin{equation}\label{eq:th-uplift}
  u(\mathbf{x}) = \Pr\left( y = 1 \mid \mathbf{x}, \text{không liên hệ} \right) - \Pr\left( y = 1 \mid \mathbf{x}, \text{có liên hệ} \right),
\end{equation}
một đại lượng chỉ ước lượng được khi có dữ liệu thí nghiệm~\parencite{ascarza2018,devriendt2021}. Olist không ghi nhận can thiệp nào, nên đề tài chỉ ước lượng được số hạng thứ nhất của~\eqref{eq:th-uplift} trong điều kiện không có hành động; giá trị của danh sách ưu tiên vì thế được đánh giá bằng Recall@$b$, và mọi diễn giải đều dừng ở mức tương quan dự báo.
